\section{Change control}

\guidance{Update the SRD when the team \emph{consciously} changes the mission
or success criteria --- not silently during CAD work.}

\subsection{When to update}

\begin{itemize}
  \item Mission goal or success criteria change
  \item Propellant, recovery concept, or major architecture change
  \item Competition rule change that affects Must requirements
  \item Approved waiver that temporarily allows non-compliance with a Must
\end{itemize}

\subsection{Process}

\begin{enumerate}[leftmargin=*]
  \item Propose change: ID(s), old text, new text, reason, impact on CONOPS / SAD / tests
  \item Review: affected department owners + safety
  \item Approve: project lead (and RSO if safety-related)
  \item Update this document + revision history; bump version
  \item Notify team (and update dependent docs)
\end{enumerate}

\subsection{Waivers}

\guidance{A \emph{waiver} is a temporary, written permission to proceed
\emph{without} fully meeting a Must requirement. It is not a silent skip.
If you permanently change the requirement, edit the requirement text via
the process above --- do not use a waiver for that.}

\noindent Fill this table only when a Must is intentionally not met for a
specific stage (example: ``MIS-004 waived for this static fire; cold-flow
runs next week''). If there are no waivers, write ``None'' in the first row
and leave the rest empty.

\noindent\begin{tabularx}{\linewidth}{@{}l l X l@{}}
\toprule
\textbf{Waiver ID} & \textbf{REQ ID} & \textbf{Reason / residual risk} & \textbf{Approver} \\
\midrule
\todo{None / W-001} & \todo{e.g.\ MIS-004} & \todo{why + what risk remains} & \todo{lead / RSO} \\
\bottomrule
\end{tabularx}
